\documentclass[10pt,a4paper]{amsart}
\usepackage{amsmath,amssymb,amsthm,graphicx}
\usepackage{verbatim}
\newenvironment{reference}{\comment}{\endcomment}
%\newenvironment{reference}{}{}
\newenvironment{slogan}{\comment}{\endcomment}
\newenvironment{history}{\comment}{\endcomment}

% For commutative diagrams we use Xy-pic
\usepackage[all]{xy}

% We use 2cell for 2-commutative diagrams.
\xyoption{2cell}
% Theorem environments.
%
\theoremstyle{plain}
\newtheorem{theorem}[subsection]{Theorem}
\newtheorem{proposition}[subsection]{Proposition}
\newtheorem{lemma}[subsection]{Lemma}

\theoremstyle{definition}
\newtheorem{definition}[subsection]{Definition}
\newtheorem{example}[subsection]{Example}
\newtheorem{exercise}[subsection]{Exercise}
\newtheorem{situation}[subsection]{Situation}

\theoremstyle{remark}
\newtheorem{remark}[subsection]{Remark}
\newtheorem{remarks}[subsection]{Remarks}

\numberwithin{equation}{subsection}

% Macros
%
\def\lim{\mathop{\mathrm{lim}}\nolimits}
\def\colim{\mathop{\mathrm{colim}}\nolimits}
\def\Spec{\mathop{\mathrm{Spec}}}
\def\Hom{\mathop{\mathrm{Hom}}\nolimits}
\def\Ext{\mathop{\mathrm{Ext}}\nolimits}
\def\SheafHom{\mathop{\mathcal{H}\!\mathit{om}}\nolimits}
\def\SheafExt{\mathop{\mathcal{E}\!\mathit{xt}}\nolimits}
\def\Sch{\mathit{Sch}}
\def\Mor{\mathop{\mathrm{Mor}}\nolimits}
\def\Ob{\mathop{\mathrm{Ob}}\nolimits}
\def\Sh{\mathop{\mathit{Sh}}\nolimits}
\def\NL{\mathop{N\!L}\nolimits}
\def\CH{\mathop{\mathrm{CH}}\nolimits}
\def\proetale{{pro\text{-}\acute{e}tale}}
\def\etale{{\acute{e}tale}}
\def\QCoh{\mathit{QCoh}}
\def\Ker{\mathop{\mathrm{Ker}}}
\def\Im{\mathop{\mathrm{Im}}}
\def\Coker{\mathop{\mathrm{Coker}}}
\def\Coim{\mathop{\mathrm{Coim}}}

% Boxtimes
%
\DeclareMathSymbol{\boxtimes}{\mathbin}{AMSa}{"02}

%
% Macros for moduli stacks/spaces
%
\def\QCohstack{\mathcal{QC}\!\mathit{oh}}
\def\Cohstack{\mathcal{C}\!\mathit{oh}}
\def\Spacesstack{\mathcal{S}\!\mathit{paces}}
\def\Quotfunctor{\mathrm{Quot}}
\def\Hilbfunctor{\mathrm{Hilb}}
\def\Curvesstack{\mathcal{C}\!\mathit{urves}}
\def\Polarizedstack{\mathcal{P}\!\mathit{olarized}}
\def\Complexesstack{\mathcal{C}\!\mathit{omplexes}}
% \Pic is the operator that assigns to X its picard group, usage \Pic(X)
% \Picardstack_{X/B} denotes the Picard stack of X over B
% \Picardfunctor_{X/B} denotes the Picard functor of X over B
\def\Pic{\mathop{\mathrm{Pic}}\nolimits}
\def\Picardstack{\mathcal{P}\!\mathit{ic}}
\def\Picardfunctor{\mathrm{Pic}}
\def\Deformationcategory{\mathcal{D}\!\mathit{ef}}
\usepackage{cleveref}
\crefname{originalsection}{original section}{original sections}
\Crefname{originalsection}{Original section}{Original sections}
\makeatletter
\newlabel{example-fibre-product-pushouts}{{05PJ}{source}{}{sourcepointer-1}{}}
\newlabel{example-fibre-product-pushouts@cref}{{[originalsection][0][]05PJ}{[1][0][]source}{}{}{}}
\makeatother
\title{A square is cartesian or cocartesian exactly when a sequence through the direct sum is exact}
\author{The Stacks Project Authors}
\date{stacks-project:homology.tex:local-snapshot-20261012; GNU FDL 1.2 or later}
\begin{document}
\maketitle
\noindent\textit{Editorial note.} All delivered mathematical text is the author's own, copied from the pinned Stacks Project chapter file homology.tex; no statement or proof step was rewritten, reordered or removed.\par\medskip
\noindent\textit{Editorial note.} The proof is the author's complete argument for the cartesian statement, including the comparison of the two component conditions; the closing clause 'and then (2) follows by duality' is the author's own routine duality remark on the part written out in full inside the counted span, and it is preserved verbatim.\par\medskip
\noindent\textit{Editorial note.} Every labelled reference occurring inside the counted spans is carried in assembly.external\_references as a non-load-bearing navigation pointer (same-chapter ones at their definition site, cross-chapter ones at the invocation inside the counted span), so no load-bearing mathematics is replaced by a stub. No bibliography citation occurs inside the counted spans.\par\medskip
\noindent\textit{Editorial note.} The same-name check for this card scanned every staged Stacks chapter of this repository -- all sources/stacks-*/pinned-original-source.tex files, the card's own chapter excluded -- and the labels found are recorded in dedup.same\_name\_cross\_chapter on the card.\par\medskip
\noindent\textit{Editorial note.} The preamble is the chapter's own native preamble excerpt already used for this source: the comment-type reference, slogan and history environments (native preamble lines 8-13), the Xy-pic setup (native preamble lines 14-18) and the native theorem declarations with the mathematical macros (native preamble lines 154-217, the actual last line of the file).\par\medskip
\noindent\textit{Allowed standard prerequisites (editorial).}\begin{itemize}
\item Fibre products and pushouts in an abelian category. Same chapter, Example example-fibre-product-pushouts. Supplies the description of the fibre product of the given square as the kernel of the map from the direct sum, which is the identification compared with exactness in the proof.
\item Direct sums and their canonical maps. Same chapter, the description of a direct sum of two objects by its two injections and two projections. Supplies the four morphisms, in particular the two projections used to compare the two isomorphisms occurring in the proof.
\end{itemize}\medskip
\ifdefined\hypertarget\hypertarget{sourcepointer-1}{}\fi
\noindent\textit{Source reference.} Original section 05PJ: sources/stacks-homology/pinned-original-source.tex, source lines 918--918. Navigation pointer invoked at proof line 1055; the counted span contains the pointer only and the cited result is not reproduced inside it, so the pointer is registered as a non-load-bearing navigation stub in the established form for Stacks material. No bibliography citation is involved.\par\medskip
\begin{lemma}
\label{lemma-characterize-cartesian}
Let $\mathcal{A}$ be an abelian category. Let
$$
\xymatrix{
w\ar[r]^f\ar[d]_g
& y\ar[d]^h\\
x\ar[r]^k
& z
}
$$
be a commutative diagram. 
\begin{enumerate}
\item The diagram is cartesian if and only if 
$$
0 \to w \xrightarrow{(g, f)} x \oplus y \xrightarrow{(k, -h)} z
$$
is exact.
\item The diagram is cocartesian if and only if 
$$
w \xrightarrow{(g, -f)} x \oplus y \xrightarrow{(k, h)} z \to 0
$$
is exact.
\end{enumerate}
\end{lemma}

\begin{proof}
Let $u = (g, f) : w \to x \oplus y$ and $v = (k, -h) : x \oplus y \to z$. 
Let $p : x \oplus y \to x$ and $q : x \oplus y \to y$ be the canonical 
projections. Let $i : \Ker(v) \to x \oplus y$ be the canonical 
injection. By Example \ref{example-fibre-product-pushouts}, the diagram is 
cartesian if and only if there exists an isomorphism 
$r : \Ker(v) \to w$ with $f \circ r = q \circ i$ and 
$g \circ r = p \circ i$. The sequence 
$0 \to w \overset{u} \to x \oplus y \overset{v} \to z$ is exact if and 
only if there exists an isomorphism $r : \Ker(v) \to w$ with 
$u \circ r = i$. But given $r : \Ker(v) \to w$, we have 
$f \circ r = q \circ i$ and $g \circ r = p \circ i$ if and 
only if $q \circ u \circ r= f \circ r = q \circ i$ and 
$p \circ u \circ r = g \circ r = p \circ i$, hence if and only if
$u \circ r = i$. This proves (1), and then (2) follows by duality. 
\end{proof}

\end{document}
